\documentclass{article}

\usepackage{iclr2027_conference,times}
\input{math_commands.tex}

\usepackage{amsmath}
\usepackage{amssymb}
\usepackage{amsthm}
\usepackage{mathtools}
\usepackage{microtype}
\usepackage{url}
\usepackage{booktabs}
\usepackage{array}
\usepackage{tabularx}
\usepackage{flafter}

\newtheorem{proposition}{Proposition}
\newtheorem{lemma}{Lemma}
\newcommand{\grain}{\textsc{Grain}}
\newcommand{\tiller}{\textsc{Tiller}}
\newcommand{\frob}[1]{\left\lVert #1\right\rVert_{\mathrm F}}
\newcommand{\inner}[2]{\left\langle #1,#2\right\rangle}

\title{Loss-Informed Transformer Training\\
with Learned Rational Activations}

\author{Anonymous authors}

\begin{document}

\maketitle

\begin{abstract}
Learning an activation changes how weight updates propagate through a network.
We use this evolving response to improve Transformer pretraining through
loss-informed coordination of matrix updates.
GRAIN (Groupwise Rational Activation with Internal Normalization) learns
rational response shapes in normalized feature coordinates and restores the
group scale at its output. TILLER (Tangent-Informed Loss Ledger for
Equal-budget Reweighting) uses these learned responses to orient matrix
directions and measure their loss contributions at shared activation positions.
A compact loss ledger retains a history of these responses, capturing how
directions reinforce or offset one another at common positions. Combining
this information with current gradients, TILLER jointly allocates update
strength across groups and layers while preserving the combined direction
norm of the feed-forward updates at each step. In language-model pretraining, GRAIN with TILLER achieves
the lowest mean validation loss among the evaluated adaptive and matrix-based
optimizer baselines across every tested corpus and model size at matched
training-token budgets.
\end{abstract}

\section{Introduction}

Learning an activation function gives a Transformer control over both its
feature representation and the response of those features to weight updates.
Fixed nonlinearities such as GELU and the SiLU gate in SwiGLU provide
effective transformations \citep{hendrycks2016gelu,shazeer2020glu}.
Rational activations make the response shape trainable
\citep{molina2020pade}, with approximation advantages that motivate their
use as expressive neural representations
\citep{boulle2020rational,tang2026expressivity}. Group-shared rational
functions bring this flexibility to Transformer feed-forward blocks
\citep{yang2025kat}. The changing response also supplies information for
optimizing the matrices around it: learning the nonlinearity changes how
an update propagates through the network.

The incoming and outgoing projections expose two complementary paths.
An incoming update changes features through the activation Jacobian;
an outgoing update recombines the activated features. Relating
these changes to the downstream loss gives a functional description of
each matrix direction. The interaction between directions is equally
important. Two groups can have the same average descent signal while
their loss responses reinforce or offset one another at shared
activation positions. Separate descent scores leave these relationships
unresolved. Measuring the responses together makes them available for
choosing the strengths of simultaneous updates.

We introduce GRAIN, Groupwise Rational Activation with Internal
Normalization, as the learned representation for this construction.
Each feature group evaluates its own rational function on
RMS-normalized coordinates and multiplies the result by the same
normalization scale \citep{zhang2019rmsnorm}. The learned coefficients
control the response to relative feature magnitudes, while the output
retains a dependence on input amplitude. Its Jacobian describes both
the local rational slope and the effect of changing the group scale.
The features and their derivatives provide the two response paths used
to guide optimization.

Matrix-aware optimization uses matrix structure and activation sensitivity
to shape updates \citep{martens2015kfac,vyas2025soap,liu2025muon}.
TILLER, Tangent-Informed Loss Ledger for Equal-budget Reweighting,
uses the learned responses to construct matrix directions and coordinate
their joint effect on the loss. It uses the response
geometry to rotate an orthogonalized momentum direction along a
factorized adaptive tangent, with orientation selected by the gradient
\citep{shazeer2018adafactor}. It then records
the selected feed-forward directions' loss responses at shared
positions in a compact temporal ledger. Together with the current
gradients, this history defines a joint model for allocating update
strength across groups and layers. The allocation preserves the
combined norm of that step's feed-forward directions, so the ledger
changes where the update acts according to the responses it produces.

This construction gives two training methods. Full TILLER continually
adapts directions and their joint allocation throughout training.
The staged method applies TILLER for an initial interval, then continues
from the learned model with a chosen optimizer. GRAIN remains trainable in both
methods. Full TILLER maintains loss-response coordination along the
complete trajectory; staged training combines a coordinated training stage
with the update rule and computational cost of its continuation
optimizer.

Across five pretraining corpora and both evaluated model sizes, GRAIN
with full TILLER achieves the lowest mean validation loss among the
tested activation--optimizer configurations at matched token budgets.
The comparisons include AdamW, Muon, SOAP, and factorized adaptive
optimizers. At the larger model size, the evaluated staged
configuration also achieves lower mean validation loss than every
baseline, while using less total training time than full TILLER on
every corpus.

[TODO: Experiment section]

\section{Related Work}

Pad\'{e} Activation Units learn rational curves jointly with network weights
\citep{molina2020pade}. Kolmogorov--Arnold Transformer (KAT) implements
group-shared rational functions followed by linear maps; its reported
configuration shares numerator coefficients across groups \citep{yang2025kat}.
GRAIN evaluates each group's rational function in RMS-normalized
coordinates \citep{zhang2019rmsnorm} and restores the group scale at its
output. It retains group sharing while learning both numerator and
denominator coefficients independently for each group.

Approximation theory establishes parameter-efficiency separations between
rational and ReLU networks \citep{boulle2020rational}.
\emph{Rational Neural Networks have Expressivity Advantages}, by Maosen Tang
and Alex Townsend, extends these comparisons to bounded-parameter
smooth-activation networks and analyzes coefficient redundancy under
learned affine normalization
\citep{tang2026expressivity}. GRAIN uses an internal normalization--rescaling
map with no additional learned affine transform.

Muon orthogonalizes momentum \citep{jordan2024muon}, while Adafactor
estimates gradient second moments from row and column statistics
\citep{shazeer2018adafactor}. TILLER combines these direction constructions
using the learned activation response. Loss-weighted features and Jacobians
determine the extent of a rotation from the polar direction along an
adaptive tangent, with its orientation selected by the current gradient.

Joint update models differ in their search directions and the information
used to compare them. K-FAC constructs parameter-space Fisher approximations from activation and
backpropagated-derivative moments, including cross-layer blocks
\citep{martens2015kfac}. Newton--Muon derives input-second-moment
preconditioning from a matrix quadratic surrogate \citep{du2026newtonmuon}.
Krylov Subspace Descent optimizes the nonlinear objective over a
curvature-generated Krylov subspace \citep{vinyals2012krylov}.
TILLER places its quadratic model on coefficients multiplying selected
pairs of incoming and outgoing group directions. Current gradients supply
the linear term. A discounted history of shared-position response probes,
augmented between probes by current-gradient information, supplies the
quadratic penalty. Its cross terms couple coefficients across groups
and layers; the allocation preserves
the combined feed-forward direction norm before weight decay.

Optimizer switching combines update rules along a training trajectory.
SWATS selects an Adam-to-SGD transition using a projection-based trigger
\citep{keskar2017switching}. Our staged method applies TILLER until a chosen
step, then continues with a selected optimizer while retaining the learned
parameters and trainable GRAIN activations.



\section{Method}
\label{sec:method}

GRAIN learns the feed-forward activation response; TILLER uses this response
to construct and coordinate matrix updates. We develop their shared
construction and its full-horizon and staged training procedures.

\subsection{GRAIN: groupwise rational activations}

GRAIN (Groupwise Rational Activation with Internal Normalization) learns
an activation in group-normalized coordinates and carries the group scale
into its output. Each feature group in a Transformer feed-forward block
has its own rational curve, following groupwise parameter sharing
\citep{yang2025kat}.

For a token representation $x\in\mathbb R^d$, partition the intermediate
features into $G$ groups of width $w$. Group $g$ has incoming and outgoing
projections $A_g\in\mathbb R^{w\times d}$ and
$B_g\in\mathbb R^{d\times w}$. Writing $z_g=A_gx$, the block computes
\begin{equation}
 \rho_g=\sqrt{\|z_g\|_2^2/w+\epsilon_{\rm act}},\qquad
 h_g=\rho_g\phi_g(z_g/\rho_g),\qquad
 y=\sum_{g=1}^G B_gh_g.
 \label{eq:overview-grain}
\end{equation}
Here $\rho_g$ is the root mean square for this token and group, stabilized
by $\epsilon_{\rm act}>0$. Normalization makes relative feature values
the input to the learned curve; multiplication by the same $\rho_g$
reintroduces the group's scale. For a linear response $\phi_g(u)=a u$,
these operations cancel exactly, giving $h_g=a z_g$.

The scalar response uses learned rational coefficients
\citep{molina2020pade} and is applied elementwise:
\begin{equation}
 \phi_g(u)=
 \frac{\sum_{j=0}^{5}a_{gj}u^j}
 {1+\sum_{j=1}^{4}|b_{gj}|\,|u|^j}.
 \label{eq:overview-rational}
\end{equation}
Its denominator is at least one, excluding real poles. Each group learns
its own numerator and denominator coefficients, shared across its features and token
positions. The functions $\phi_g$ start from a common SiLU approximation
and are trained jointly with the projection weights.

The RMS depends on all features in a group, so an incoming update can
couple their responses through the full group Jacobian
$J_g=\partial h_g/\partial z_g$. An outgoing update recombines the
current features $h_g$. TILLER uses the Jacobian and features to
orient matrix updates and evaluate their first-order effects on the loss.
Appendix~\ref{sec:grain} derives the Jacobian and its action on a direction.

\subsection{TILLER: loss-informed update coordination}

TILLER uses the learned activation response to orient matrix directions
and choose their strengths jointly across groups and layers.

\subsubsection{Response-guided matrix directions}

For $L$ layers, index the $C=LG$ layer--group pairs by $i=(\ell,g)$, with incoming and
outgoing roles $r\in\{A,B\}$. For each matrix, TILLER forms a polar
direction $p_i^r$ and a factorized adaptive direction $a_i^r$ from
the same momentum \citep{jordan2024muon,liu2025muon,shazeer2018adafactor}.
The polar map rebalances singular values; factorized adaptation rescales
entries using row and column gradient second moments. Removing the
adaptive direction's component parallel to the polar direction gives
a tangent along which to rotate the update.

Two loss-weighted response statistics in $[0,1]$ control the rotation.
Normalized participation $\pi_i^r$ measures the breadth of response energy: across Jacobian modes
for the incoming role and feature coordinates for the outgoing role.
Congruence $\kappa_i^r$ is a weighted Frobenius cosine between the current response
kernels, $J_iJ_i^\top$ or $h_ih_i^\top$, and those evaluated with the
initial rational coefficients at the same current inputs. Suppressing
the role index, the rotation candidate is
\begin{equation}
 \begin{gathered}
 a_i^\perp=a_i-\frac{\inner{a_i}{p_i}}{\frob{p_i}^2}p_i,
 \qquad e_i=\pi_i(1-\kappa_i^2),\\
 \widetilde q_i=\sqrt{1-e_i}\,p_i+
 \sigma_i\sqrt{e_i}\frac{\frob{p_i}}{\frob{a_i^\perp}}a_i^\perp .
 \end{gathered}
 \label{eq:overview-rotation}
\end{equation}
Here $\sigma_i\in\{-1,1\}$ makes the tangent's contraction with the
current gradient nonnegative. Orthogonality preserves the parent norm.
An unchanged response gives $e_i=0$; a broad response with low congruence
permits a larger rotation. A finite, nondegenerate candidate is selected
when both it and its parent have positive gradient contraction; otherwise
the parent is retained. Positive norm calibration gives the direction
$q_i$ used below. Appendix~\ref{sec:directions} specifies these operations.

\subsubsection{Joint loss-response allocation}

One signed coefficient $c_i$ scales both directions $(q_i^A,q_i^B)$
of each group, preserving the relative calibration of its two paths.
To measure interactions among groups, sample $K$ example--token
positions shared across layers. At position $k$, let $x_{k\ell}$ be
the block input and $\xi_{k\ell}$ the derivative of the summed minibatch
loss with respect to the block output, scaled by the gradient-clipping
factor. The response matrix $S_t\in\mathbb R^{K\times C}$ contains
\begin{equation}
 S_{t,ki}=\inner{\xi_{k\ell}}
 {B_iJ_{ki}q_i^Ax_{k\ell}+q_i^Bh_{ki}}.
 \label{eq:overview-score}
\end{equation}
Each row collects the groups' contributions through one activation
position to the downstream loss derivative. For coefficients $c$, $S_tc$ combines these
contributions. The cross terms in $\|S_tc\|_2^2/K$ distinguish
reinforcement and cancellation that separate mean descent scores omit.

The loss ledger retains a discounted estimate $G_t$ of the response
second moment $S_t^\top S_t/K$ using a Robust Frequent Directions sketch
\citep{ghashami2016frequent,luo2019robust}.
Its columns keep the same layer--group identities as directions evolve.
Probes refreshed every eight steps measure positional variation; synthetic rows aligned with
current gradient contractions update the history between probes, retaining
the latest measured response scale. Together with the
current descent signal
$s_{t,i}=\inner{g_i^A}{q_i^A}+\inner{g_i^B}{q_i^B}$,
this history defines the coefficient model
\begin{equation}
 \begin{gathered}
 F_t(c)=-\eta s_t^\top c+
 \frac{\eta^2}{2}\bigl(c^\top G_tc+2r_t^\top c\bigr),\\
 \sum_i\omega_ic_i^2=\sum_i\omega_i,\qquad
 \omega_i=\frob{q_i^A}^2+\frob{q_i^B}^2 .
 \end{gathered}
 \label{eq:overview-coordinate-model}
\end{equation}
Here $g_i^A,g_i^B$ are clipped gradients, $\eta$ is the learning rate,
and $r_t$ records the response cross term with weight decay at rate
$\lambda_{\rm wd}$.
The linear term rewards current descent. The quadratic term penalizes
historical squared responses of the combined directions and weight decay.
Its cross terms make the allocation sensitive to where recorded directions
reinforce or cancel one another. The constraint preserves the combined norm of the paired
feed-forward directions before weight decay. A projected solve proposes
coefficients; acceptance requires a lower $F_t$ than the all-ones
allocation, the norm constraint, and positive layerwise gradient and
momentum contractions for both allocations. Otherwise the all-ones
allocation is retained. Appendix~\ref{sec:loss-measurements} gives the
ledger and solver.

The selected coefficients update each incoming projection as
$A_i\leftarrow(1-\eta\lambda_{\rm wd})A_i-\eta c_iq_i^A$,
and analogously for $B_i$. Attention uses response-guided rotations
with layerwise GRAIN summaries. Rational coefficients and the remaining
parameter classes use AdamW \citep{loshchilov2019decoupled}.

Full TILLER applies these updates throughout the training horizon $T$.
Its directions and joint allocation combine periodically refreshed
GRAIN responses with current gradient information at every step,
maintaining loss-response coordination along the full trajectory.

Staged training applies the same updates through a chosen step $T_s<T$,
then continues the feed-forward and attention matrices with a selected
optimizer. It retains the learned model and global learning-rate
schedule, with an optimizer-specific state transition. GRAIN remains
trainable, and its AdamW state continues across the transition. This
method combines a coordinated training stage with the update rule and
computational cost of its continuation optimizer.
Table~\ref{tab:algorithm} presents both procedures.



\section{Experiments}

[TODO: Experiment section]

\section{Discussion}

\section{Conclusion}

\subsection*{AI use statement}

Generative AI tools assisted with method development, experimental planning,
implementation, mathematical derivations and proof checking, literature
analysis, result interpretation, and manuscript drafting and revision.

\bibliography{references}
\bibliographystyle{iclr2027_conference}

\clearpage
\appendix
\section{GRAIN response}
\label{sec:grain}
\label{app:grain-calculus}

The normalization and rescaling in GRAIN let the learned curve operate
on relative feature magnitudes while retaining a response to input
amplitude. For one group, suppress the layer and group indices and write
\begin{equation}
 \Phi(z)=\rho\phi(u),\qquad
 u=z/\rho,\qquad
 \rho=(z^\top z/w+\epsilon_{\rm act})^{1/2}.
 \label{eq:grain-normalize}
\end{equation}
The scalar function acts coordinatewise, and its coefficients are held
fixed when differentiating with respect to $z$. Its denominator is
positive for finite coefficients. Absolute-value derivatives at zero
use $\operatorname{sign}(0)=0$; at these kinks the Jacobian below is
the selected autodifferentiation linearization.

Both the normalized argument and the restored scale contribute to the
response. With $f=\phi(u)$, $d_\phi=\phi'(u)$, and
$r_\phi=f-u\odot d_\phi$, differentiation gives
\begin{equation}
 \nabla_z\rho=u/w,\qquad
 \frac{\partial u}{\partial z}=\rho^{-1}(I-uu^\top/w),\qquad
 J=\frac{\partial\Phi}{\partial z}
 =\operatorname{Diag}(d_\phi)+\frac{r_\phi u^\top}{w}.
 \label{eq:grain-jacobian}
\end{equation}
The diagonal term applies the local rational slopes. The rank-one term
couples the coordinates through the shared scale, with $r_\phi$
measuring the departure from a local linear response.

TILLER needs the action of this response on perturbations and loss
sensitivities. These actions follow directly from the rank-one form:
\begin{equation}
 Jv=d_\phi\odot v+r_\phi\frac{u^\top v}{w},\qquad
 J^\top\zeta=d_\phi\odot\zeta+
 u\frac{r_\phi^\top\zeta}{w}.
 \label{eq:grain-jacobian-actions}
\end{equation}
Each needs one inner product and coordinatewise products after the
curve and slope have been evaluated. The incoming path in
$B\Phi(Ax)$ consequently acts as $BJUx$ for a matrix direction $U$,
while an outgoing direction $V$ acts as $V\Phi(Ax)$.

The coefficient initialization in the algorithm approximates SiLU on
$[-5,5]$ and is shared only at initialization: all ten coefficients
subsequently learn independently in each group. The implemented width
is $w=256$. For a SwiGLU intermediate width $H$, the GRAIN width is
$H_r=Gw=3H/2$, so its two dense projections contain
$2dH_r=3dH$ weights, matching the three SwiGLU projections.

\section{Response-guided matrix directions}
\label{sec:directions}

TILLER chooses a matrix direction before assigning its joint update
coefficient. The direction combines two transformations of the same
momentum: approximate polar normalization controls its singular values,
and factorized adaptation incorporates the history of entrywise gradient
scales. The learned activation determines how much of the adaptive
component enters; the current gradient determines its orientation.
This construction preserves the calibrated parent norm, leaving the
coefficient allocation to redistribute that norm across groups.

For a layer--group pair $i=(\ell,g)$ and role $r\in\{A,B\}$, let
$g_{t,i}^r$ denote the accumulated, distributed-averaged, clipped
gradient. We use $\beta_1=\beta_2=0.95$ and
$\epsilon_{\rm opt}=10^{-8}$; $\delta_{32}$ denotes the smallest
positive normal float32 value. From a zero buffer, form
\begin{equation}
 v_{t,i}^r=\beta_1v_{t-1,i}^r+(1-\beta_1)g_{t,i}^r,
 \qquad
 m_{t,i}^r=\beta_1v_{t,i}^r+(1-\beta_1)g_{t,i}^r.
 \label{eq:nesterov-momentum}
\end{equation}
The parent is a groupwise Muon direction
\citep{jordan2024muon,liu2025muon},
\begin{equation}
 p_{t,i}^r=\chi_{\rm grp}\mathcal P_5(m_{t,i}^r),\qquad
 \chi_{\rm grp}=\sqrt{\frac{\min\{H_r,d\}}{G\min\{w,d\}}},\qquad
 \mu_{\rm mlp}=0.2\sqrt{\max\{H_r,d\}},
 \label{eq:group-polar-calibration}
\end{equation}
where $\mathcal P_5$ is the five-step Newton--Schulz map in the
algorithm and $H_r=Gw$. With exact full-rank polar factors, the sum
of squared block norms would be $G\min\{w,d\}$; multiplying by
$\chi_{\rm grp}$ matches the nominal whole-projection norm
$\min\{H_r,d\}$. The multiplier $\mu_{\rm mlp}$ is applied after
direction selection.

Factorized second moments provide the adaptive transformation
\citep{shazeer2018adafactor}. Suppressing the matrix indices and
initializing both marginals to zero, set
\begin{equation}
 \begin{gathered}
 R_t=\beta_2R_{t-1}+(1-\beta_2)\operatorname{rowsum}(g_t^2),\qquad
 C_t=\beta_2C_{t-1}+(1-\beta_2)\operatorname{colsum}(g_t^2),\\
 V_{t,jk}=\frac{R_{t,j}C_{t,k}}
 {\max\{\sum_sR_{t,s},\delta_{32}\}(1-\beta_2^t)},
 \qquad a_t=m_t/(\sqrt{V_t}+\epsilon_{\rm opt}).
 \end{gathered}
 \label{eq:factored-direction}
\end{equation}
The marginals commute with exponential averaging, so they summarize the
corresponding row and column sums of the full second-moment history.
Their normalized outer product supplies the entrywise scale. Removing
$a_t$'s component parallel
to $p_t$ then gives a tangent to the Frobenius sphere through $p_t$.

The response statistics describe where that tangent should contribute.
At a sampled position, incoming perturbations pass through $J_{ki}$,
whereas outgoing perturbations recombine $h_{ki}$. We quantify their
response breadth, for nonzero responses, by
\begin{equation}
 \operatorname{pr}_A(J)=
 \frac{\operatorname{tr}(JJ^\top)^2}{w\|JJ^\top\|_{\rm F}^2},\qquad
 \operatorname{pr}_B(h)=\frac{\|h\|_2^4}{w\|h\|_4^4}.
 \label{eq:response-summary-components}
\end{equation}
The first measures energy spread across squared singular values, and
the second across squared feature coordinates. Equal energy in $m$
nonzero components gives $m/w$. Thus the incoming statistic is spectral,
while the outgoing statistic refers to the model's feature coordinates.

Positions are weighted by the loss sensitivity of the corresponding
path. For the downstream derivative $\xi_{k\ell}$, define
\begin{equation}
 \alpha^A_{ki}=\|J_{ki}^\top B_i^\top\xi_{k\ell}\|_2^2/w,
 \qquad \alpha^B_{ki}=\|\xi_{k\ell}\|_2^2/d,
 \qquad
 \pi_i^r=\frac{\sum_k\alpha^r_{ki}\operatorname{pr}_r(\mathcal R^r_{ki})}
 {\sum_k\alpha^r_{ki}},
 \label{eq:participation-main}
\end{equation}
where $\mathcal R^A_{ki}=J_{ki}$ and $\mathcal R^B_{ki}=h_{ki}$,
and the total sensitivity in each ratio is positive.
The incoming weight is the mean squared derivative at the preactivation;
the outgoing weight is the corresponding quantity at the block output.
This weighting emphasizes response geometry at loss-sensitive positions.

To measure how learning changes the response, evaluate the initial
rational coefficients at the same current normalized inputs and scales.
For $\Gamma^A_{ki}=J_{ki}J_{ki}^\top$ and
$\Gamma^B_{ki}=h_{ki}h_{ki}^\top$, let the superscript $0$ denote
this reference evaluation. For positive weighted kernel norms, define
congruence by
\begin{equation}
 \kappa_i^r=
 \frac{\sum_k\alpha^r_{ki}\langle\Gamma^r_{ki},\Gamma^{r,0}_{ki}\rangle}
 {\sqrt{\left(\sum_k\alpha^r_{ki}\|\Gamma^r_{ki}\|_{\rm F}^2\right)
 \left(\sum_k\alpha^r_{ki}\|\Gamma^{r,0}_{ki}\|_{\rm F}^2\right)}}.
 \label{eq:congruence-main}
\end{equation}
This is a cosine between collections of response kernels weighted by
$\sqrt{\alpha^r_{ki}}$. Evaluating both curves at current inputs isolates
the coefficient change.

Set $e_i^r=\pi_i^r(1-(\kappa_i^r)^2)$. For a nonzero parent and
tangent, suppressing time and role indices, the direction before
calibration is
\begin{equation}
 a_i^\perp=a_i-\frac{\langle a_i,p_i\rangle}{\|p_i\|_{\rm F}^2}p_i,
 \qquad
 \widetilde q_i=\sqrt{1-e_i}\,p_i+
 \sigma_i\sqrt{e_i}\frac{\|p_i\|_{\rm F}}{\|a_i^\perp\|_{\rm F}}a_i^\perp.
 \label{eq:tangent-chord}
\end{equation}
Set $\sigma_i=1$ when $\langle g_i,a_i^\perp\rangle\geq0$
and $\sigma_i=-1$ otherwise. The numerical algorithm computes the combination
as coefficients of $p_i$ and $a_i$, and selects it when it is finite,
nondegenerate, and has a positive gradient contraction along with the
parent. Calibration then gives $q_i=\mu_{\rm mlp}\widetilde q_i$;
a rejected rotation uses $q_i=\mu_{\rm mlp}p_i$.

Attention uses the response of the same layer's GRAIN map as a shared
signal. It averages participation across groups but aggregates the kernel
inner products and squared norms before taking congruence. With these
layer statistics, its rotation fraction is
\begin{equation}
 e_\ell^{\rm att}=
 \sqrt{\overline\pi_\ell^A\overline\pi_\ell^B}
 (1-\kappa_\ell^A\kappa_\ell^B).
 \label{eq:attention-response}
\end{equation}
The breadth factor includes both feed-forward paths, and departure in
either can increase the rotation. QKV and attention-output matrices
share this fraction while retaining their own moments, tangents, and
gradient-selected signs. Their selected directions are applied directly;
joint coefficients are reserved for the paired GRAIN projections.

\section{Joint loss-response coordination}
\label{sec:loss-measurements}

One coefficient per layer--group pair controls the strength of both its
incoming and outgoing directions. Directions in disjoint parameter
blocks can nevertheless act together through the network. Measuring
their contributions at shared activation positions makes those
interactions available to the allocation model, while using the full
gradient supplies its current descent term.

\subsection{Shared-position measurements}

Fix a group $i=(\ell,g)$, a position $k$, and directions $(q_i^A,q_i^B)$.
Differentiating
$(B_i+\alpha q_i^B)\Phi_i((A_i+\alpha q_i^A)x_{k\ell})$
at $\alpha=0$ gives
\begin{equation}
 \begin{aligned}
 \delta y_{ki}&=B_iJ_{ki}q_i^Ax_{k\ell}+q_i^Bh_{ki},\\
 S_{t,ki}&=\langle q_i^Ax_{k\ell},J_{ki}^\top B_i^\top\xi_{k\ell}\rangle
 +\langle q_i^Bh_{ki},\xi_{k\ell}\rangle.
 \end{aligned}
 \label{eq:functional-score}
\end{equation}
The cotangent $\xi_{k\ell}$ is the derivative of the summed minibatch
loss at the block output, multiplied by the realized clipping factor.
It collects the hidden position's influence on downstream loss terms.
The score consequently measures a contribution through that activation
position along the paired direction.

Probe derivatives and parameter gradients must use consistent scaling.
For a microbatch with $N$ flattened positions whose mean loss is divided
by the accumulation count $M$, a captured derivative $\bar\xi_{k\ell}$
contains the factor $1/(NM)$. TILLER uses
\begin{equation}
 \xi_{k\ell}=NM\chi_{\rm clip}\bar\xi_{k\ell},\qquad
 \chi_{\rm clip}=\min\{1,(\|g_{\rm raw}\|_2+10^{-6})^{-1}\}.
 \label{eq:cotangent-rescaling}
\end{equation}
The clipping factor is held fixed while evaluating directions. Averaging
the resulting scores over all positions and processes recovers the
clipped minibatch derivative under this averaging convention. Its exact
value for the current directions is already available from the parameter
gradients:
\begin{equation}
 s_{t,i}=\langle g_{t,i}^A,q_{t,i}^A\rangle
 +\langle g_{t,i}^B,q_{t,i}^B\rangle.
 \label{eq:current-score}
\end{equation}
The linear term uses $s_t$ directly; the $K$ sampled rows of $S_t$
measure how the positional contributions vary together.

Weight decay contributes through the same response paths. Substituting
$(\lambda_{\rm wd}A_i,\lambda_{\rm wd}B_i)$ for $(q_i^A,q_i^B)$
and summing over groups and layers gives the decay response
$a_t^{\rm dec}\in\mathbb R^K$. The combined squared response is
\begin{equation}
 \frac{\|S_tc+a_t^{\rm dec}\|_2^2}{K}
 =c^\top\frac{S_t^\top S_t}{K}c
 +2c^\top\frac{S_t^\top a_t^{\rm dec}}{K}
 +\frac{\|a_t^{\rm dec}\|_2^2}{K}.
 \label{eq:functional-response-energy}
\end{equation}
The last term is independent of the allocation. The first two therefore
supply a quadratic response penalty and its decay cross term. Combined
with current descent and a fixed direction norm, this penalty compares
the recorded joint effects of different allocations at the scale of
the applied step.

\subsection{Discounted response history}
\label{sec:ledger}

Responses are measured at steps $1,9,17,\ldots$, using $K=32$
positions aligned across layers. A refresh recomputes participation,
congruence, and positional scores at the current parameters. Between
refreshes, only the response statistics are reused: the matrix moments,
selected directions, and full-gradient contractions are recomputed at
every step. Historical columns preserve their layer--group identities,
with each row recording the direction in use when it was measured.

The intervening steps update the ledger's orientation using current
descent while retaining the latest measured response scale. Let
$\nu=\|S_{\rm ref}\|_{\rm F}/\sqrt K$ and let $d_t^{\rm dec}$ be
the clipped-gradient contraction with the GRAIN decay direction. The
synthetic rows are
\begin{equation}
 S_t=\boldsymbol1_K(\gamma_ts_t)^\top,\qquad
 a_t^{\rm dec}=\boldsymbol1_K\gamma_td_t^{\rm dec},\qquad
 \gamma_t=\begin{cases}\nu/\|s_t\|_2,&\|s_t\|_2>0,\\
 0,&\|s_t\|_2=0.
 \end{cases}
 \label{eq:trace-matched-rows}
\end{equation}
For nonzero $s_t$, their second moment has rank at most one and trace
$\nu^2$. The refreshed rows supply positional variation; the synthetic
rows carry current-gradient information between those measurements.
The numerical zero threshold is specified in the algorithm.

For the resulting sequence of measured and synthetic rows, the
uncompressed target is
\begin{equation}
 G_t^\star=\beta_2G_{t-1}^\star+(1-\beta_2)S_t^\top S_t/K,
 \qquad G_1^\star=S_1^\top S_1/K.
 \label{eq:ledger-target}
\end{equation}
Thus $c^\top G_t^\star c$ averages squared responses of a coefficient
combination over the recorded history. The decay $\beta_2=0.95$
discounts old relationships as the directions and learned responses
change.

The ledger represents the previous metric by a row factor and an
isotropic remainder,
$\widehat G_{t-1}=\mathcal B_{t-1}^\top\mathcal B_{t-1}/K+
\tau_{t-1}I$. Adding current rows gives the metric used for selection:
\begin{equation}
 \begin{aligned}
 \mathcal A_t&=\begin{bmatrix}\sqrt{\beta_2}\mathcal B_{t-1}\\
 \sqrt{1-\beta_2}S_t\end{bmatrix},&
 G_t&=\mathcal A_t^\top\mathcal A_t/K+\beta_2\tau_{t-1}I,\\
 r_t&=\beta_2r_{t-1}+(1-\beta_2)S_t^\top a_t^{\rm dec}/K.
 \end{aligned}
 \label{eq:current-ledger-metric}
\end{equation}
At the first step, $\mathcal A_1=\mathcal B_1=S_1$,
$\tau_0=\tau_1=0$, and $r_1=S_1^\top a_1^{\rm dec}/K$.
The current metric $G_t$ contains the fresh rows before their
compression; $\widehat G_t$ denotes the compressed metric retained
for future steps.

Robust Frequent Directions limits the stored factor to $r=64$ rows
\citep{ghashami2016frequent,luo2019robust}. Let the singular values
of $\mathcal A_t$ be $\sigma_1\geq\sigma_2\geq\cdots$.
Let $r_*$ be the smaller of $r$ and the number of singular values.
When the rank exceeds $r$, set $\Delta_t=\sigma_{r+1}^2$; otherwise
set $\Delta_t=0$. Store
\begin{equation}
 \mathcal B_t=\operatorname{Diag}
 \bigl(\sqrt{\sigma_1^2-\Delta_t},\ldots,
 \sqrt{\sigma_{r_*}^2-\Delta_t}\bigr)V_{1:r_*}^\top,
 \qquad
 \tau_t=\beta_2\tau_{t-1}+\Delta_t/(2K).
 \label{eq:rfd-reconstruction}
\end{equation}
For rank at most $r$, this retains an exact factor.
The numerical algorithm obtains the factor from the row Gram matrix.
Its dominant rows retain response combinations, while the isotropic
term centers the approximation error across coefficient directions.

\subsection{Fixed-norm allocation}
\label{sec:transaction}

The allocation compares different strengths at the norm of the current
unit-allocation update. With
$\omega_i=\|q_i^A\|_{\rm F}^2+\|q_i^B\|_{\rm F}^2$ and
$\mathcal E_0=\sum_i\omega_i$, disjoint parameter blocks give
$\|\Delta\theta(c)\|_2^2=\eta_t^2\sum_i\omega_ic_i^2$.
Consequently, the constraint $\sum_i\omega_ic_i^2=\mathcal E_0$
matches the feed-forward direction norm of $c=\boldsymbol1$
before weight decay. Signed coefficients let the model compare the
combined actions of groups and layers while preserving the relative
calibration of each group's incoming and outgoing paths.

Whitening the constraint turns the proposal calculation into a spherical
quadratic problem. For positive $\omega_i$, set
$W=\operatorname{Diag}(\omega)$ and $y=W^{1/2}c$. Dividing the
coefficient model by $\eta_t^2$ gives
\begin{equation}
 \begin{gathered}
 \mathcal Q_t(y)=\tfrac12y^\top M_ty-b_t^\top y,
 \qquad \|y\|_2^2=\mathcal E_0,\\
 M_t=W^{-1/2}G_tW^{-1/2},\qquad
 b_t=W^{-1/2}(s_t/\eta_t-r_t).
 \end{gathered}
 \label{eq:whitened-transaction}
\end{equation}
Its stationarity equation is $(M_t+\lambda I)y=b_t$.
The positive-multiplier branch gives a positive-definite system because
$M_t$ is positive semidefinite. TILLER uses this branch to construct
a proposal in a subspace of at most $\min\{32,C\}$ dimensions.
The proposal is defined for positive group energies, positive learning
rate, and nonzero $b_t$.

The subspace is built from repeated metric actions starting at $b_t$
using two reorthogonalization passes and orthogonal restarts
\citep{lanczos1950iteration}. The factor representation gives
\begin{equation}
 M_tv=W^{-1/2}\mathcal A_t^\top
 (\mathcal A_tW^{-1/2}v)/K+\beta_2\tau_{t-1}W^{-1}v.
 \label{eq:ledger-metric-action}
\end{equation}
Let $Q_k$ be the resulting basis and $T_k$ the accumulated symmetric
tridiagonal matrix. Clipping
its negative eigenvalues gives
$T_k^+=U\operatorname{Diag}(\theta)U^\top$ with $\theta_j\geq0$.
For $\gamma=U^\top Q_k^\top b_t$, the projected proposal is
\begin{equation}
 y=Q_kU\operatorname{Diag}((\theta+\lambda)^{-1})\gamma,
 \qquad
 \psi(\lambda)=\sum_j\frac{\gamma_j^2}{(\theta_j+\lambda)^2}
 =\mathcal E_0.
 \label{eq:secular-equation}
\end{equation}
The prescribed bracket and $64$ bisection steps produce
$c=W^{-1/2}y$ when a positive root is bracketed. If the lower-end test
rules out a root in this bracket, TILLER retains the unit allocation.

The projected model proposes coefficients; the original factor metric
decides whether to use them. Acceptance requires a finite candidate,
$F_t(c)<F_t(\boldsymbol1)$, and scaled norm residual at most
$10^{-8}$, measured in float64. For every layer and role, it also requires
\begin{equation}
 \sum_g c_{\ell g}\langle g_{\ell g}^r,q_{\ell g}^r\rangle>0,
 \qquad
 \sum_g c_{\ell g}\langle m_{\ell g}^r,q_{\ell g}^r\rangle>0,
 \qquad r\in\{A,B\}.
 \label{eq:transaction-contractions}
\end{equation}
The same inequalities must hold for $c=\boldsymbol1$.
Rejected proposals retain the unit allocation; accepted coefficients
are cast to float32 before application to the matrix directions.

\section{Properties of the construction}
\label{app:analysis}

\subsection{Activation and update geometry}

The GRAIN Jacobian separates changes in relative features from changes
in group amplitude. A direction $v$ with $u^\top v=0$ leaves the scale
unchanged to first order and satisfies $Jv=d_\phi\odot v$. Along the input,
$u^\top u/w=1-\epsilon_{\rm act}/\rho^2$ yields
\begin{equation}
 Jz=\Phi(z)-\frac{\epsilon_{\rm act}}{\rho}r_\phi.
 \label{eq:grain-radial}
\end{equation}
For $\epsilon_{\rm act}=0$ and nonzero $z$, positive rescaling obeys
$\Phi(az)=a\Phi(z)$ for $a>0$, and the identity becomes
$Jz=\Phi(z)$. The correction term quantifies the stabilizer's effect
on the radial response. For a linear curve $\phi(u)=au$,
$r_\phi=0$ and $\Phi(z)=az$, so the scale operations cancel even
with a positive stabilizer.

Participation and congruence define a bounded rotation fraction. For the
ratios in equation~\ref{eq:response-summary-components}, nonnegative
component energies $v_j$ satisfy
$\sum_jv_j^2\leq(\sum_jv_j)^2\leq w\sum_jv_j^2$, placing
nonzero participation in $[1/w,1]$. The identities
$\langle JJ^\top,J^0J^{0\top}\rangle=\|J^\top J^0\|_{\rm F}^2$
and $\langle hh^\top,h^0h^{0\top}\rangle=(h^\top h^0)^2$,
together with Cauchy--Schwarz, place defined congruence in $[0,1]$.
The zero-response conventions in the algorithm keep $e_i^r$ in this
interval and retain the parent when reference alignment is undefined.
\label{app:response-proof}

Congruence measures alignment of response energy. Sign reversals of $J$
or $h$ preserve their kernels, and a common nonzero gain
$\phi=c\phi^0$ gives $\kappa=1$. The loss-response measurements
retain the gain through $J$ and $h$.

Orthogonality gives $\|\widetilde q_i\|_{\rm F}^2=
(1-e_i)\|p_i\|_{\rm F}^2+e_i\|p_i\|_{\rm F}^2$.
Choosing $\sigma_i$ by the sign of $\langle g_i,a_i^\perp\rangle$
makes the tangent contribute nonnegatively to the descent contraction
for a subtraction update. The calibrated selection has norm
$\|\mu_{\rm mlp}p_i\|_{\rm F}$ in exact arithmetic when numerical
floors are inactive; parent fallback has that norm directly.

Participation attenuates the angular departure from the reference. With
$e_i^r=\pi_i^r(1-(\kappa_i^r)^2)$, the rotation angle $\vartheta_i^r$
and kernel angle $\psi_i^r=\arccos\kappa_i^r$ obey
$\sin^2\vartheta_i^r=\pi_i^r\sin^2\psi_i^r$.
A concentrated response therefore permits a smaller adaptive component,
while full participation permits $\vartheta_i^r=\psi_i^r$.

The coefficient proposal enforces its prescribed norm through the scalar
equation $\psi(\lambda)=\mathcal E_0$. For nonzero $\gamma$ and
$\lambda>0$,
\begin{equation}
 \psi'(\lambda)=-2\sum_j\frac{\gamma_j^2}{(\theta_j+\lambda)^3}<0.
 \label{eq:secular-monotonicity}
\end{equation}
Thus the positive branch has at most one root. For an orthonormal
basis and nonnegative projected eigenvalues, the algorithm's upper
endpoint satisfies
$\psi(\lambda_{\rm hi})\leq
\|b_t\|_2^2/\lambda_{\rm hi}^2\leq\mathcal E_0$.
The lower-end test establishes whether this bracket contains a root.

\subsection{Response-history approximation}

The retained factor approximates a specified discounted response history.
Its error can be bounded independently of coefficient selection because
Robust Frequent Directions applies the same singular-value shrinkage
to each retained direction and stores half that amount isotropically.
The following result concerns this representation in exact arithmetic.

\begin{proposition}[Response-metric approximation]
\label{prop:ledger-error}
For $K>0$, $\beta_2\in[0,1]$, and $\tau_{t-1}\geq0$, let
$G_t$ and $(\mathcal B_t,\tau_t)$ be defined by
equations~\ref{eq:current-ledger-metric} and~\ref{eq:rfd-reconstruction}.
Then $\widehat G_t=\mathcal B_t^\top\mathcal B_t/K+\tau_tI$ is
positive semidefinite, and for every $c\in\mathbb R^C$,
\begin{equation}
 \|\widehat G_t-G_t\|_2\leq\frac{\Delta_t}{2K},\qquad
 |c^\top(\widehat G_t-G_t)c|
 \leq\frac{\Delta_t}{2K}\|c\|_2^2.
 \label{eq:rfd-error}
\end{equation}
\end{proposition}
\begin{proof}
Both terms of $\widehat G_t$ are positive semidefinite. For
$\Delta_t=0$, the two metrics coincide. For positive $\Delta_t$,
work in a right-singular-vector basis of $\mathcal A_t$, completed
by its null space. Each retained squared singular value loses
$\Delta_t/K$ and receives $\Delta_t/(2K)$ from the isotropic
term, giving error $-\Delta_t/(2K)$. A discarded direction with
squared singular value $s\in[0,\Delta_t]$ has error
$(\Delta_t/2-s)/K$; a null direction has error $\Delta_t/(2K)$.
Every absolute eigenvalue of the difference is therefore at most
$\Delta_t/(2K)$. The quadratic-form bound follows from
$|c^\top Ec|\leq\|E\|_2\|c\|_2^2$ for symmetric $E$.
\end{proof}

The isotropic increment centers the one-step compression error. The
same retained factor with the previous remainder
$\beta_2\tau_{t-1}I$ has error eigenvalues in
$[-\Delta_t/K,0]$. Adding $\Delta_t I/(2K)$ centers this interval
at zero and halves the one-compression spectral bound, from
$\Delta_t/K$ to $\Delta_t/(2K)$, for the response metric retained
for subsequent steps.


Discounting limits how older compression errors accumulate. Write
$E_t=\widehat G_t-G_t^\star$ and
$D_t=\widehat G_t-G_t$. The common initialization gives $E_1=0$,
and the two recurrences imply $E_t=\beta_2E_{t-1}+D_t$. Hence
\begin{equation}
 \|E_t\|_2\leq
 \sum_{j=2}^t\beta_2^{t-j}\frac{\Delta_j}{2K},
 \qquad G_t-G_t^\star=\beta_2E_{t-1}.
 \label{eq:rfd-history-error}
\end{equation}
The current metric carries only previous compression error because it
includes fresh rows before the new shrinkage. These bounds describe
fidelity to the measured-and-synthetic response history and its
quadratic penalty.

\section{Complete training algorithm}
\label{app:implementation}
\label{app:algorithm}
\label{app:numerical}

Full-horizon and staged training use the same TILLER update. A horizon
$T$ and transition index $T_s\in\{1,\ldots,T\}$ specify how long it is
applied. When $T_s<T$, a continuation rule $\mathcal U$ and its state
initializer $\mathcal I$ complete the procedure. The initializer maps the
parameters and optimizer state after update $T_s$ to the moments,
counters, and matrix grouping needed by $\mathcal U$. The model and
global learning-rate schedule continue across this boundary, as does
the AdamW state for rational coefficients and the remaining parameter classes.

Rows 1--7 of Table~\ref{tab:algorithm} give the training loop;
the remaining rows define its local operations.

\begingroup
\small
\setcounter{LTchunksize}{100}
\setlength{\tabcolsep}{4pt}
\begin{longtable}{@{}r >{\raggedright\arraybackslash}p{0.91\linewidth}@{}}
\caption{GRAIN training with full-horizon or staged TILLER.
All local operations used by the training loop are specified in this
table. Setting $T_s=T$ gives full-horizon TILLER; otherwise
$(\mathcal U,\mathcal I)$ defines the continuation.}
\label{tab:algorithm}\\
\toprule
& Procedure and local operations \\
\midrule
\endfirsthead
\multicolumn{2}{l}{\tablename\ \thetable\ (continued)}\\
\toprule
& Procedure and local operations \\
\midrule
\endhead
\bottomrule
\endfoot
& Inputs: initial model parameters, $L$ layers, $G$ groups of width
$w=256$, model width $d$, $H_r=Gw$, $C=LG\geq2$, horizon $T$,
transition $T_s\in\{1,\ldots,T\}$, schedule $\eta_t$, matrix decay $\lambda_{\rm wd}$,
a model and loss averaged over the microbatch's $N$ token positions,
a minibatch iterator, $M=4$ accumulation steps, $P$ distributed processes,
and per-parameter AdamW schedules $\eta_\theta(t)$ and decay rates.
The grouped matrices have $A_i\in\mathbb R^{w\times d}$ and
$B_i\in\mathbb R^{d\times w}$, $i=(\ell,g)$.  For $T_s<T$,
supply $\mathcal U$ and $\mathcal I$. Require $\eta_t>0$ during
TILLER. The supported attention layout has stacked QKV of shape
$3d\times d$ and an output matrix of shape $d\times d$, with 16 heads,
four heads per polar block, and $16\mid d$.
\\[2pt]
& Constants: $\beta_1=\beta_2=0.95$, $\epsilon_{\rm act}=10^{-6}$,
$\epsilon_{\rm opt}=10^{-8}$, $K=32$, refresh interval $8$, persistent
row limit $64$, Krylov width $\min\{32,C\}$, and $64$ bisection steps.
The matrix decay default is $\lambda_{\rm wd}=0.1$. Rational
coefficients, tied embeddings, normalization parameters, and biases
have zero AdamW decay; other AdamW matrices use $0.1$.
Write $\delta_{32}$ and $\delta_{64}$ for the smallest positive normal
float32 and float64 values, and $\varepsilon_{64}$ for float64 machine
epsilon. Inner products of matrices are Frobenius inner products.
Elementwise powers, products, and division are used for moment updates.
\\[2pt]
& For every group, initialize the numerator and denominator coefficients,
in ascending powers, to
\[
\begin{aligned}
a^0={}&(2.4915714620520875\!\times\!10^{-5},\ 0.5000668663485823,\\
&0.2499441908995771,\ 0.0526200910151016,\\
&0.005525765640115002,\ 0.00024199321178747251),\\
b^0={}&(0.00022554017910995938,\ 0.10511420350691994,\\
&2.8656992394109812\!\times\!10^{-5},\ 0.0004816962110562025).
\end{aligned}
\]
During the forward pass, each group computes $z_g=A_gx$,
$\rho_g=\sqrt{\|z_g\|^2/w+\epsilon_{\rm act}}$,
$h_g=\rho_g P_g(z_g/\rho_g)/D_g(z_g/\rho_g)$ with the polynomials
defined in row 9, and the block returns $\sum_gB_gh_g$.
Keep a frozen copy $(a^0,b^0)$ for response comparisons. Initialize each
matrix's momentum $v$ and row/column second moments $R,C_{\mathrm{col}}$ to
zero, initialize AdamW moments and counts to zero, and leave the response
ledger and response cache empty. \\
\midrule
1 & For $t=1,\ldots,T$, clear parameter gradients and let the refresh flag be
$t\leq T_s$ and $(t-1)\bmod8=0$. At the current parameters,
evaluate each microbatch mean loss,
divide by the accumulation count $M$, and backpropagate. On a refresh,
capture the probe positions and their block inputs, preactivations,
features, and block-output derivatives using the selection rule below.
Distributed-average the accumulated parameter gradients. Evaluate all
responses and directions before applying any parameter update. \\
2 & With $g_{\rm raw}$ the concatenated gradient, set
$\chi=\min\{1,(\|g_{\rm raw}\|_2+10^{-6})^{-1}\}$ and
$g=\chi g_{\rm raw}$. Nonfinite gradients are errors. For a captured
microbatch with $N$ flattened positions and captured derivative
$\bar\xi$, set $\xi=NM\chi\bar\xi$. \\
3 & If $t\leq T_s$, refresh and cache the response summaries in rows 8--12
when the refresh flag is set. Construct all feed-forward and attention
directions using rows 13--15. \\
4 & When $t\leq T_s$, for each $i=(\ell,g)$ and $r\in\{A,B\}$, form
$d_i^r=\langle g_i^r,q_i^r\rangle$,
$u_i^r=\langle m_i^r,q_i^r\rangle$,
$s_i=d_i^A+d_i^B$ (summing the two float32 contractions in float64), and
$\omega_i=\|q_i^A\|_{\rm F}^2+\|q_i^B\|_{\rm F}^2$.
Construct current response rows and the ledger factor using rows 16--19. \\
5 & Still within $t\leq T_s$, apply rows 20--23 to the current factor,
isotropic term, decay cross term, $(d_i^r,u_i^r)$, $\omega$, and
$\eta_t$, obtaining $c$. Retain $(\mathcal B,\tau,r)$ for the next
step. Apply
\[
A_i\leftarrow(1-\eta_t\lambda_{\rm wd})A_i-\eta_tc_iq_i^A,
\quad
B_i\leftarrow(1-\eta_t\lambda_{\rm wd})B_i-\eta_tc_iq_i^B.
\]
For each attention matrix $W$, apply
$W\leftarrow(1-\eta_t\lambda_{\rm wd})W-\eta_tq_W$.
Directions and selected coefficients are float32. Retain the unscaled
selected direction $\widetilde q$ and its matrix-size multiplier $\mu$
from row 15. For application, cast $c_i\widetilde q_i^r$ (or
$\widetilde q_W$ for attention) to the parameter dtype, then add it
with scalar multiplier $-\eta_t\mu$. \\
6 & If $t>T_s$, update the feed-forward and attention matrices with
the supplied rule $\mathcal U$, its state, and the current clipped
gradients at learning rate $\eta_t$. In either branch, update rational coefficients and the
remaining parameters by the AdamW operation below. \\
7 & After completing update $t=T_s<T$, set the continuation
state to $\mathcal I(\text{current parameters},\text{current optimizer state})$,
retaining the learned parameters and AdamW
state. Stop response collection and ledger updates. Continue the loop
with the same global schedule. After $t=T$, return the trained model.
Rows 8--23 and the AdamW operation define the local operations invoked
by rows 1--7. \\
\midrule
& \emph{Probe selection and response summaries.} \\*
8 & For $P$ processes, process $p\in\{0,\ldots,P-1\}$ contributes
$K_p=\lfloor(p+1)K/P\rfloor-\lfloor pK/P\rfloor$ positions.
Require each selection to satisfy $0\leq m\leq N$.
From each of the $M$ microbatches capture $m_p=\lceil K_p/M\rceil$
positions, then select $K_p$ from their concatenation. Selecting $m>1$
rows from $N$ rows uses zero-based indices
$\lfloor j(N-1)/(m-1)\rfloor$, $j=0,\ldots,m-1$; $m=1$ selects
zero and $m=0$ selects nothing. Use identical indices across layers
and retain the captured tensors locally. \\
9 & Cast captured records and rational coefficients to float32; evaluate
the response and statistics in rows 9--12 in float32. At each selected
position and group, set
$\rho=\sqrt{z^\top z/w+\epsilon_{\rm act}}$ and $u=z/\rho$.
For current or frozen coefficients evaluate
\[
\begin{gathered}
P(u)=\sum_{j=0}^{5}a_ju^j,\qquad
D(u)=1+\sum_{j=1}^{4}|b_j||u|^j,\\
f=P(u)/D(u),\quad
d_\phi=\frac{P'(u)D(u)-P(u)D'(u)}{D(u)^2},\quad
r_\phi=f-u\odot d_\phi,\\
P'(u)=\sum_{j=1}^{5}ja_ju^{j-1},\qquad
D'(u)=\sum_{j=1}^{4}j|b_j||u|^{j-1}\operatorname{sign}(u).
\end{gathered}
\]
Use $\operatorname{sign}(0)=0$. For response statistics set $h^{\rm eval}=\rho f$ and
$J=\operatorname{Diag}(d_\phi)+r_\phi u^\top/w$.
Apply $Jv=d_\phi\odot v+r_\phi(u^\top v)/w$ and
$J^\top v=d_\phi\odot v+u(r_\phi^\top v)/w$ without storing $J$.
Frozen evaluation uses the same current $u$ and $\rho$.
Retain the captured forward feature $h$ separately for row 16. \\
10 & Define $\Gamma^A=JJ^\top$ and $\Gamma^B=h^{\rm eval}(h^{\rm eval})^\top$ through their
inner products. The diagonal-plus-rank-one Jacobian makes $JJ^\top$ a
diagonal matrix plus a rank-two correction, so these inner products
reduce to vector contractions and a $2\times2$ matrix calculation.
For two Jacobians at the same $u$, let
$r=r_\phi/w$, $\widetilde r=\widetilde r_\phi/w$,
$v=d_\phi\odot u$, $\widetilde v=\widetilde d_\phi\odot u$,
$s=u^\top u$, $Z=[v,r]^\top[\widetilde v,\widetilde r]$, and
$E=\left[\begin{smallmatrix}0&1\\1&s\end{smallmatrix}\right]$.
Then evaluate
\[
\begin{aligned}
\langle\Gamma^A,\widetilde\Gamma^A\rangle={}&
\langle d_\phi^2,\widetilde d_\phi^2\rangle
+2\langle d_\phi^2,\widetilde v\odot\widetilde r\rangle
+s\langle d_\phi^2,\widetilde r^2\rangle\\
&+2\langle\widetilde d_\phi^2,v\odot r\rangle
+s\langle\widetilde d_\phi^2,r^2\rangle
+\langle EZE,Z\rangle,\\
\operatorname{tr}(JJ^\top)={}&\|d_\phi\|^2+2\langle d_\phi\odot u,r\rangle
+s\|r\|^2,\\
\langle\Gamma^B,\widetilde\Gamma^B\rangle={}&((h^{\rm eval})^\top\widetilde h^{\rm eval})^2.
\end{aligned}
\]
Self-inner products give the squared kernel norms. \\
11 & For each role compute per-position participation and sensitivity
\[
\begin{gathered}
P^A_{ki}=\frac{\operatorname{tr}(J_{ki}J_{ki}^\top)^2}
 {w\max\{\|J_{ki}J_{ki}^\top\|_{\rm F}^2,\delta_{32}\}},\quad
P^B_{ki}=\frac{\|f_{ki}\|_2^4}{w\max\{\|f_{ki}\|_4^4,\delta_{32}\}},\\
\alpha^A_{ki}=\|J_{ki}^\top B_i^\top\xi_{k\ell}\|^2/w,
\qquad \alpha^B_{ki}=\|\xi_{k\ell}\|^2/d.
\end{gathered}
\]
Clip only the participation ratios to $[0,1]$; a zero response gives zero. Here $f=h^{\rm eval}/\rho$,
which gives the same outgoing participation as $h^{\rm eval}$ when the floor is
inactive. Accumulate
$D_i^r=\sum_k\alpha^r_{ki}$,
$N_i^r=\sum_k\alpha^r_{ki}P^r_{ki}$,
$X_i^r=\sum_k\alpha^r_{ki}\langle\Gamma^r_{ki},\Gamma^{r,0}_{ki}\rangle$,
$Y_i^r=\sum_k\alpha^r_{ki}\|\Gamma^r_{ki}\|_{\rm F}^2$,
$Z_i^r=\sum_k\alpha^r_{ki}\|\Gamma^{r,0}_{ki}\|_{\rm F}^2$.
Sum these statistics across processes. \\
12 & Set $\pi_i^r=\operatorname{clip}_{[0,1]}(N_i^r/\max(D_i^r,\delta_{32}))$.
If either role has $D_i^r=0$, set both $\pi_i^r$ to one.
Set $\kappa_i^r=\operatorname{clip}_{[0,1]}
(X_i^r/\max(\sqrt{Y_i^rZ_i^r},\delta_{32}))$.
If either unfloored denominator $\sqrt{Y_i^rZ_i^r}$ is zero, or the group's coefficients
equal their initialization, set both congruences to one.
For layer $\ell$, set
$\overline\pi_\ell^r=G^{-1}\sum_g\pi_{\ell g}^r$ and
$\kappa_\ell^r=(\sum_gX_{\ell g}^r)/
\max\{\sqrt{(\sum_gY_{\ell g}^r)(\sum_gZ_{\ell g}^r)},\delta_{32}\}$,
clipped to $[0,1]$. A zero unfloored denominator
$\sqrt{(\sum_gY_{\ell g}^r)(\sum_gZ_{\ell g}^r)}$ for either role sets both
layer congruences to one; so does equality to initialization for every
group. Cache
\[
e_i^r=\pi_i^r(1-(\kappa_i^r)^2),\qquad
e_\ell^{\rm att}=\sqrt{\overline\pi_\ell^A\overline\pi_\ell^B}
(1-\kappa_\ell^A\kappa_\ell^B).
\]
All response statistics must be finite. \\
\midrule
& \emph{Matrix moments and direction construction.} \\*
13 & For each TILLER-assigned matrix, update its first moment in the gradient dtype:
$v\leftarrow\beta_1v+(1-\beta_1)g$ and
$m=\beta_1v+(1-\beta_1)g$.
For each feed-forward group block, and for each whole attention matrix,
update
\[
\begin{gathered}
R\leftarrow\beta_2R+(1-\beta_2)\operatorname{rowsum}(g^2),\quad
C_{\mathrm{col}}\leftarrow\beta_2C_{\mathrm{col}}+(1-\beta_2)\operatorname{colsum}(g^2),\\
V=\frac{RC_{\mathrm{col}}^{\top}}
{\max\{\boldsymbol1^\top R,\delta_{32}\}(1-\beta_2^t)},
\qquad a=m/(\sqrt V+\epsilon_{\rm opt}).
\end{gathered}
\]
Second moments and the direction arithmetic use float32. \\
14 & To compute $\mathcal P_5(M)$, transpose a tall block $M$,
cast to bfloat16, compute its Frobenius norm in float32, floor the norm
at $10^{-7}$, cast the norm to bfloat16, and divide the block by it.
Starting with the resulting $X$, perform five bfloat16 iterations
\[
H=XX^\top,\qquad X\leftarrow3.4445X+
(-4.775H+2.0315H^2)X.
\]
Restore the original orientation. For each GRAIN matrix block set
$p=\chi_{\rm grp}\mathcal P_5(m)$ with
$\chi_{\rm grp}=\sqrt{\min\{H_r,d\}/(G\min\{w,d\})}$.
For QKV, apply $\mathcal P_5$ to four contiguous row blocks of width
$d/4$ in each of Q, K, and V, reassemble, and multiply by $1/\sqrt3$.
For the attention output matrix, apply it to four contiguous column
blocks of width $d/4$ and reassemble. Apply these parent calibration
factors in bfloat16, then cast the parents to float32. \\
15 & Use $e=e_i^r$ for feed-forward block $(i,r)$ and
$e=e_\ell^{\rm att}$ for both attention matrices in layer $\ell$.
For the corresponding parent $p$, adaptive direction $a$, current
gradient $g$, and cached $e\in[0,1]$, compute
\[
\begin{gathered}
P=\|p\|_{\rm F}^2,\quad A=\|a\|_{\rm F}^2,\quad
h=\langle p,a\rangle,\quad \zeta=h/\max(P,\delta_{32}),\\
D_\perp=\max\{A-h^2/\max(P,\delta_{32}),0\},\quad
\sigma=\begin{cases}1,&\langle g,a\rangle-\zeta\langle g,p\rangle\geq0,\\-1,&\text{otherwise},\end{cases}\\
c_a=\sigma\sqrt e\sqrt{P/\max(D_\perp,\delta_{32})},\quad
c_p=\sqrt{1-e}-c_a\zeta,\quad
D_{\mathrm{cand}}=c_p\langle g,p\rangle+c_a\langle g,a\rangle.
\end{gathered}
\]
Use $(c_p,c_a)$ only if $p,m,g,a,c_p,c_a,D_{\mathrm{cand}}$ are finite,
$P,A>0$, $D_\perp>\delta_{32}$,
$\langle g,p\rangle>0$, and $D_{\mathrm{cand}}>0$; otherwise use $(1,0)$.
The selected direction $\widetilde q=c_pp+c_aa$ must be finite. Set
$\mu=0.2\sqrt{\max\{H_r,d\}}$ for GRAIN blocks and
$\mu=0.2\sqrt{\max\{n,m\}}$ for a whole $n\times m$ attention matrix.
The calibrated direction is $q=\mu\widetilde q$. \\
\midrule
& \emph{Response rows and temporal ledger.} \\*
16 & On refresh steps, for every locally selected position and group form
\[
\begin{aligned}
S_{ki}&=\langle q_i^Ax_{k\ell},J_{ki}^\top B_i^\top\xi_{k\ell}\rangle
+\langle q_i^Bh_{ki},\xi_{k\ell}\rangle,\\
a_k^{\rm dec}&=\lambda_{\rm wd}\sum_i
\bigl(\langle A_ix_{k\ell},J_{ki}^\top B_i^\top\xi_{k\ell}\rangle
+\langle B_ih_{ki},\xi_{k\ell}\rangle\bigr).
\end{aligned}
\]
Gather $S$ and $a^{\rm dec}$ across processes to obtain the $K$ aligned
rows. Save $\nu=\|S\|_{\rm F}/\sqrt K$, accumulating its squared norm in
float64 and storing the result in float32. \\
17 & Between refreshes compute in float32
$d^{\rm dec}=\lambda_{\rm wd}\sum_i
(\langle g_i^A,A_i\rangle+\langle g_i^B,B_i\rangle)$.
In float64 compute $\gamma=\nu/\|s\|_2$ if
$\|s\|_2>\delta_{64}$ and $\gamma=0$ otherwise. Cast
$\gamma s$ and $\gamma d^{\rm dec}$ to float32 and repeat them
in all $K$ rows:
$S=\boldsymbol1_K(\gamma s)^\top$,
$a^{\rm dec}=\boldsymbol1_K\gamma d^{\rm dec}$.
Require finite rows, contractions, and $\nu\geq0$. \\
18 & At the first step set $\mathcal A=S$, $\delta=0$,
$r=S^\top a^{\rm dec}/K$, and save
$(\mathcal B,\tau)=(S,0)$.
On later steps, with previous ledger $(\mathcal B_-,\tau_-,r_-)$,
form
\[
\mathcal A=\begin{bmatrix}\sqrt{\beta_2}\mathcal B_-\\
\sqrt{1-\beta_2}S\end{bmatrix},\qquad
\delta=\beta_2\tau_-,\qquad
r=\beta_2r_-+(1-\beta_2)S^\top a^{\rm dec}/K.
\]
The current metric is $G=\mathcal A^\top\mathcal A/K+\delta I$.
Keep $\mathcal A$ in float32; accumulate the decay cross term in
float64 and store it in float32. Require finite factors and cross terms
and nonnegative $\delta$. Keep $\delta$ in float64 for row 19 and pass
a float32 copy, together with the stored $r$, to the current coefficient
operation. \\
19 & After the first step, form the next-step factor from the float64 symmetric Gram matrix
$H=(\mathcal A\mathcal A^\top+
(\mathcal A\mathcal A^\top)^\top)/2$ and diagonalize it as
$H=U\operatorname{Diag}(\lambda)U^\top$, with eigenvalues sorted in
increasing order and clipped below at zero. Let $n$ be its dimension,
$j=\min\{64,n\}$, and $\Delta=\lambda_{n-j}$ when $n>j$,
otherwise $\Delta=0$ (indices start at one). For the largest $j$
eigenvalues $\lambda_+$ and corresponding columns $U_+$, save
\[
\mathcal B=
\operatorname{Diag}\!\left(\sqrt{
\frac{\max(\lambda_+-\Delta,0)}{\max(\lambda_+,\delta_{64})}}
\right)U_+^\top\mathcal A,\qquad
\tau=\delta+\Delta/(2K).
\]
Compute the diagonal
$\operatorname{diag}(\mathcal B^\top\mathcal B/K)+\tau\boldsymbol1$
in float64, then cast it and $\mathcal B,\tau$ to float32 for storage.
Require finite output and $\tau,\Delta\geq0$.
The first step retains $(S,0)$ directly. \\
\midrule
& \emph{Projected coefficient proposal and acceptance.} \\*
20 & Require finite $\mathcal A,\delta,r,d_i^r,u_i^r,\omega$, with
$\delta\geq0$, every $\omega_i>0$, and $\eta>0$.
Convert inputs to float64. Set
$d_\omega=\sqrt{\omega}$, $E_0=\sum_i\omega_i$,
$b=(s/\eta-r)\oslash d_\omega$, and
$\mathcal Z_{:i}=\mathcal A_{:i}/(\sqrt K\,d_{\omega,i})$,
where $\oslash$ denotes elementwise division.
Require finite $\|b\|_2>0$.
Invalid inputs are errors. Use the matrix action
\[
\mathcal M(v)=\mathcal Z^\top(\mathcal Z v)
+(\delta\oslash\omega)\odot v.
\]
\\
21 & Let $k=\min\{32,C\}$, $v_1=b/\max(\|b\|_2,\delta_{64})$,
$v_0=0$, $\beta_0=0$, and $Q_0$ have no columns.
For $j=1,\ldots,k$, set $Q_j=[Q_{j-1},v_j]$, compute $a_j=\mathcal M(v_j)$,
$\alpha_j=v_j^\top a_j$, and
$z=a_j-\alpha_jv_j-\beta_{j-1}v_{j-1}$.
Perform $z\leftarrow z-Q_j(Q_j^\top z)$ twice and set
$\widehat\beta_j=\|z\|_2$.
If $j<k$, set
$\epsilon_j=64\varepsilon_{64}\max\{\|a_j\|_2,1\}$ and
start with $v_{j+1}=z/\max(\widehat\beta_j,\delta_{64})$.
For $j<k$, when $\widehat\beta_j\leq\epsilon_j$, take the seed
$f_a=\cos(aj)$ for $a=1,\ldots,C$, perform
$f\leftarrow f-Q_j(Q_j^\top f)$ twice, and replace $v_{j+1}$ by
$f/\|f\|_2$ if $\|f\|_2>\epsilon_j$.
For $j<k$, set $\beta_j=0$ when
$\widehat\beta_j\leq\epsilon_j$, and
$\beta_j=\widehat\beta_j$ otherwise. \\
22 & Form symmetric tridiagonal $T_k$ with diagonal $\alpha_j$ and
off-diagonal $\beta_j$, and $Q=Q_k$. Require finite $Q,T_k$ and
$\|Q^\top Q-I\|_{\rm F}\leq2\!\times\!10^{-10}k$; a failed
basis check is an error. Diagonalize $T_k=U\operatorname{Diag}(\theta)U^\top$,
replace $\theta$ by $\max(\theta,0)$, and set $\gamma=U^\top Q^\top b$.
Define
\[
\psi(\lambda)=\sum_j\frac{\gamma_j^2}{(\theta_j+\lambda)^2},\qquad
\lambda_{\rm lo}=\varepsilon_{64}C,\qquad
\lambda_{\rm hi}=\|b\|_2/\sqrt{E_0}+\max_j\theta_j+\lambda_{\rm lo}.
\]
Record whether $\psi(\lambda_{\rm lo})<E_0$. \\
23 & Repeat $64$ times: set
$\lambda=(\lambda_{\rm lo}+\lambda_{\rm hi})/2$; if
$\psi(\lambda)>E_0$, replace $\lambda_{\rm lo}$ by $\lambda$;
otherwise replace $\lambda_{\rm hi}$ by $\lambda$.
Construct
$c_*=[QU(\gamma/(\theta+\lambda_{\rm hi}))]\oslash d_\omega$.
Evaluate both $c_*$ and $\boldsymbol1$ in
\[
F(c)=-\eta\sum_i s_ic_i+
\frac{\eta^2}{2}
\left(\|\mathcal A c\|_2^2/K+\delta\|c\|_2^2+2r^\top c\right).
\]
Accept $c_*$ only if the recorded lower-end condition was false,
the candidate and objective are finite,
$F(c_*)<F(\boldsymbol1)$, and
$|\sum_i\omega_ic_{*,i}^2-E_0|/\max\{E_0,1\}\leq10^{-8}$.
For every layer $\ell$ and role $r$, also require
$\sum_g c_{*,\ell g}d_{\ell g}^r>0$,
$\sum_g c_{*,\ell g}u_{\ell g}^r>0$,
$\sum_g d_{\ell g}^r>0$, and $\sum_g u_{\ell g}^r>0$.
Return the accepted coefficient vector cast to float32; otherwise
return $\boldsymbol1$. \\
\midrule
& \emph{AdamW operation.} For each assigned parameter $\theta$ with
gradient $g_\theta$, current scheduled learning rate $\eta_\theta(t)$, and decay
$\lambda_\theta$, advance its count $n$ and moments
$m_\theta\leftarrow0.9m_\theta+0.1g_\theta$,
$v_\theta\leftarrow0.95v_\theta+0.05g_\theta^2$.
Apply
\[
\theta\leftarrow(1-\eta_\theta(t)\lambda_\theta)\theta-
\eta_\theta(t)\frac{m_\theta/(1-0.9^n)}
{\sqrt{v_\theta/(1-0.95^n)}+10^{-8}}.
\]
These moments and counts continue across the staged transition.
\\
\end{longtable}
\endgroup


\end{document}
